\documentclass[11pt,a4paper]{article}
\usepackage[ngerman]{babel}
\usepackage[TS1,T1]{fontenc}
\usepackage[utf8x]{inputenc}
\selectlanguage{ngerman}
\usepackage{theorem}
\usepackage[scaled=0.9]{helvet}
\usepackage{amsmath}
\usepackage{amssymb}
\usepackage{hyperref}
\usepackage{stmaryrd}
\usepackage{pgf,tikz}
\usepackage{relsize}
\usepackage{enumitem}
\usepackage{graphicx}
\usepackage{algpseudocode,amsmath,xifthen}

\newcounter{numb}
\theoremstyle{break}
\theorembodyfont{\upshape}
   	\newtheorem{exercise}{Exercise}[numb]
\setlength{\oddsidemargin}{0cm}
\setlength{\textwidth}{16cm}
\setlength{\textheight}{23cm}
\setlength{\topmargin}{-2cm}

\usetikzlibrary{shapes,arrows,automata,positioning,decorations.fractals}
\renewcommand\familydefault{\sfdefault}

\newcommand{\header}[2]{
\begin{minipage}[b]{\textwidth}
		\parbox[t]{5cm}{%
			\includegraphics[width=4cm]{unilogo}
			\hfill
		}
		\parbox[b]{11cm}{%
			%\scshape%
			Heinrich-Heine-University D\"usseldorf\\
			Computer Science Department\\
			Software Engineering and Programming Languages\\
			%Professor Dr.\ M.\ Leuschel
		Philipp K\"orner
}
		
\end{minipage}
	\begin{center}
		\bf
		Functional Programming -- WT 2024 / 2025\\
        Reading Guide {\thenumb}: {#1}
	\end{center}

    \pagestyle{empty}
    \paragraph{Timeline:} This unit should be completed by {#2}.
}

\setcounter{numb}{11}


\begin{document}

\begin{minipage}[b]{\textwidth}
	\parbox[t]{5cm}{%
		\includegraphics[width=4cm]{unilogo}
		\hfill
	}
	\parbox[b]{11cm}{%
		%\scshape%
		Heinrich-Heine-University D\"usseldorf\\
		Computer Science Department\\
		Software Engineering and Programming Languages\\
		%Professor Dr.\ M.\ Leuschel
		Philipp K\"orner
	}
\end{minipage}
\begin{center}
	\bf
	Functional Programming -- WT 2022 / 2023\\
	Reading Guide 11: The Expression Problem
\end{center}

\pagestyle{empty}

\paragraph{Timeline:} This unit should be completed by 09.01.2023.

\paragraph{Note:} This unit is way smaller than usual. Happy holidays and enjoy your time off!

\section{Material} 

\begin{itemize}
%\item Clojure for the Brave and True, chapter 13
%\item Clojure Reference: Multimethods \url{https://clojure.org/reference/multimethods}
%\item Clojure Reference: Protocols \url{https://clojure.org/reference/protocols}
\item Philip Wadler: The Expression Problem \url{http://homepages.inf.ed.ac.uk/wadler/papers/expression/expression.txt}
\item Slides: vl-10.pdf 
\end{itemize}


\section{Learning Outcomes}

After completing this unit you should be able to

\begin{itemize}
    \item describe the expression problem.
    \item identify use cases as well as advantages and disadvantages of wrappers and monkey patching.
    %\item explain how multimethods and protocols solve the expression problem.
    %\item use multimethods and protocols.
    %\item compare the expressiveness of multimethods and protocols.
    %\item recall and compare the performance of multimethods and protocols.
\end{itemize}

\section{Highlights}

\begin{itemize}
    \item Expression problem
    %\item Multimethods
    %\item Protocols
    %\item Functions / Macros: \verb|defmulti|, \verb|defmethod|, \verb|defprotocol|, \verb|extend-protocol|, \verb|extend-type|, \verb|extend|
\end{itemize}



\section{Exercises}

No practical exercises here.
Proceed with the usual 4clojure exercises
only if you find the time.


%\begin{exercise}[Lazy Searching - 4clojure Nr. 108]
%
%Implement a function that receives any number of ascending sequences
%as input and returns the smallest element that occurs in all sequences.
%The inputs can be infinite sequences.
%
%\begin{verbatim}
%(common-min [3 4 5])
%=> 3
%(common-min [1 2 3 4 5 6 7] [0.5 3/2 4 19])
%=> 4
%(common-min (range) (range 0 100 7/6) [2 3 5 7 11 13])
%=> 7
%(common-min (map (fn [x] (* x x x)) (range)) ;; Cubic numbers
%            (filter (fn [x] (zero? (bit-and x (dec x)))) (range)) ;; Powers of two
%            (iterate inc 20)) ;; Numbers greater than/equal to 20
%=> 64
%\end{verbatim}

%\end{exercise}

	\section*{Questions}
	If you have any questions, please contact Philipp K"orner (\texttt{p.koerner@hhu.de}) or post it to the Rocket.Chat group.
\end{document}

